\documentclass[11pt]{article}
\usepackage[a4paper, top=18mm, bottom=18mm, left=20mm, right=20mm]{geometry}
\usepackage[T2A]{fontenc}
\usepackage[utf8]{inputenc}
\usepackage[ukrainian]{babel}
\usepackage{graphicx}
\usepackage[export]{adjustbox}
\usepackage{amsmath}
\usepackage{amssymb}
\setlength{\parindent}{0pt}
\setlength{\parskip}{4pt}
\pagestyle{empty}
\begin{document}
%begin
{\large\textbf{Контрольна робота}}

\textbf{Варіант \No\,2}\hfill \textit{ПІБ:}\,\underline{\hspace{60mm}}
\vspace{4mm}

%beginitem
1. 

print('R', end=' ')
f=['0','1','2','3','4']
print(f.pop(0), end=' ')
print(f)


%enditem

%beginitem
2. Що буде виведено за виконання фрагменту коду?

%code
#include <iostream>

int h(int a, int b){
    int c = 84;
    if (b != -3)
        c = 5;
    else if (a > -2)
         return 2;
    else 
        return 0;
    return c;
}

int main(){
    std::cout << h(3, 3);
    return 0;
}
%code

%enditem

%beginitem
3. 
try:
    res = "False"!=True
    if isinstance(res, bool):
        print(f'bool;{res}')
    elif isinstance(res, int):
        print(f'int;{res}')
    elif isinstance(res, float):
        print(f'float;{res:.2f}')
    else:
        print('???')
except:
    print('Z')

%enditem

%beginitem
4. Що буде виведено за виконання фрагменту коду?

%code
print('R', end=' ')
f=[0,1,2,3,4,5,6,7,8,9]
print(f[:-3:-3])
%code


%enditem

%beginitem
5. %goodpath:Scripts/2019/Mod2019_1/Lex/03-good.txt
Відмітити все, що є коректним оператором С++:
( 1 ) x not is y

( 2 ) print(end="1","qwe")

( 3 ) 3**-1

( 4 ) x-=2

%enditem

%beginitem
6. 
Написати програму, що запитує у користувача значення дійсних параметрів $x$ та $y$, обчислює для них значення виразу та повiдомляє введенi данi та результат обчислення.
$$\frac{\log_{2} x - 44}{(\tg x - 0.3) (y - 99)}$$ 


%enditem

%beginitem
7. Що буде виведено за виконання фрагменту коду?

%code
f = ('a','b','c','d','e',0,1,2,3)
print(f[-1])
%code


%enditem

%beginitem
8. %goodpath:Scripts/2018/Mod2018_1/03-good.txt
Відмітити все, що є коректним оператором С++:
( 1 ) double n.m;

( 2 ) if a<b a=0;

( 3 ) {int x,y;}

( 4 ) double(2)/87

%enditem

%beginitem
9. 
a = 5
b = 2
c = 0

def h(b):
global c    a *= 3
    b = 4
    c = 2
    return a + b + c

a = 9
b = 6
c = 8

try:
    print(h(b), a, b, c, sep=';')
except:
    print('Z', a, b, c, sep=';')

%enditem

%beginitem
10. %goodpath:Scripts/2018/Mod2018_1/01-good.txt
Відмітити все, що є однією правильною лексемою:
( 1 ) 5_w

( 2 ) (1)

( 3 ) )

( 4 ) 0.5

%enditem

%end
\newpage
\end{document}

